\chapter*{Preface: How to Read This if You Are Tired and Blaming Yourself}\label{ch:preface}
\addcontentsline{toc}{chapter}{Preface: How to Read This if You Are Tired and Blaming Yourself}
\noindent This book does not ask you to optimize your life, buy a wearable, or embrace the doctrine of ``perfect sleep hygiene.''
Instead, it tries to make sense of why sleep and dreams exist at all, and why modern work, stimulants, neurodivergence, and
perpetual connectivity make an ancient biological process so fragile. If you are picking this up at 2~a.m. in a haze of caffeine
and guilt, here is the contract we are offering:

\begin{enumerate}
    \item \textbf{No moralizing.} Sleep debt is not a personal failure. It is the predictable outcome of a nervous system evolved
    for cyclic rest trying to survive in a culture of perpetual availability.
    \item \textbf{The goal is understanding, not judgment.} Throughout Parts~I--III we treat sleep as a computationally necessary
    algorithm that alternates between maintenance, consolidation, and creative generalization. Knowing the mechanics makes it
    easier to separate structural barriers (shift work, disability, caretaking) from choices that are actually yours to change.
    \item \textbf{You can skim.} Each chapter begins with a short thesis statement so that an exhausted reader can capture the
    argument in a page or two before diving deeper on a better day. Sidebars flag especially practical material, and summaries at
    the end of each part restate the key findings without academic hedging.
\end{enumerate}

If you take nothing else from the next several hundred pages, let it be this: sleep is not an optional luxury you keep failing to
``earn.'' It is a phase of thinking that minds require to stay flexible. You are already doing something heroic by trying to
understand it before collapsing.
